\documentclass[10pt]{article}
\usepackage[T1]{fontenc}
\usepackage[utf8]{inputenc}
\usepackage[french]{babel}
\usepackage{amsmath,amssymb,mathrsfs,mathpazo}
\usepackage{geometry}
\geometry{a4paper,margin=2cm}
\usepackage{graphicx,array,tabularx,tabularray,colortbl,multirow,multicol,booktabs,enumitem,pifont,ulem,fancybox,diagbox,cancel}
\usepackage{pstricks,pstricks-add,pst-plot,pst-tree,pst-node,pst-text,pst-func,pst-math,pst-bspline,pst-3dplot,pst-eucl,pst-solides3d}
\let\pscalculate\undefined
\usepackage{pst-bezier}
\usepackage{tikz,pgfplots,tkz-tab}
\usetikzlibrary{arrows,arrows.meta,patterns,calc,positioning,trees,decorations.markings,shapes.geometric,shapes.misc,shapes.arrows,intersections,decorations.pathmorphing,decorations.pathreplacing,backgrounds,shadows,babel}
\pgfplotsset{compat=1.18}
\usepgfplotslibrary{fillbetween}
\usepackage[np]{numprint}
\usepackage{listings,makecell,pas-tableur,twoopt}
\let\barre\undefined
\usepackage[pstricks]{tablvar}
\definecolor{darkgreen}{RGB}{0,100,0}
\providecommand{\red}{\color{red}}
\providecommand{\blue}{\color{blue}}
\providecommand{\green}{\color{green}}
\providecommand{\black}{\color{black}}
\providecommand{\white}{\color{white}}
\providecommand{\gray}{\color{gray}}
\providecommand{\cyan}{\color{cyan}}
\providecommand{\magenta}{\color{magenta}}
\providecommand{\yellow}{\color{yellow}}
\providecommand{\violet}{\color{violet}}
\providecommand{\orange}{\color{orange}}
\providecommand{\euro}{\texteuro}
\providecommand{\EUR}{\texteuro}
\providecommand{\ieme}{\textsuperscript{e}}
\providecommand{\degres}{\ensuremath{^\circ}}
\providecommand{\degree}{\ensuremath{^\circ}}
\providecommand{\vv}[1]{\overrightarrow{#1}}
\providecommand{\Coord}[1]{#1}
\providecommand{\decosix}{}
\providecommand{\starredbullet}{\ensuremath{\bullet}}
\pagestyle{empty}
\setlength{\parindent}{0pt}
\usepackage{adjustbox}\DeclareUnicodeCharacter{2236}{:}\DeclareUnicodeCharacter{2190}{\ensuremath{\leftarrow}}
\providecommand{\R}{\mathbb{R}}
\providecommand{\N}{\mathbb{N}}
\providecommand{\Z}{\mathbb{Z}}
\providecommand{\Q}{\mathbb{Q}}
\providecommand{\D}{\mathbb{D}}
\providecommand{\C}{\mathbb{C}}
\providecommand{\e}{{\mathrm{e}}}
\providecommand{\vect}[1]{\overrightarrow{#1}}
\providecommand{\vectt}[1]{\overrightarrow{\mathrm{#1}}}
\providecommand{\vvt}[1]{\overrightarrow{\mathrm{#1}}}
\providecommand{\np}[2][]{{#2~#1}}
\providecommand{\vv}[1]{{\overrightarrow{#1}}}
\renewcommand{\d}{{\mathrm{d}}}
\providecommand{\degres}{{^\circ}}
\providecommand{\degre}{{^\circ}}
\providecommand{\blue}{\color{blue}}
\providecommand{\red}{\color{red}}
\providecommand{\enskip}{\quad}
\providecommand{\llbracket}{\lbrack\!\lbrack}
\providecommand{\rrbracket}{\rbrack\!\rbrack}
\providecommand{\euro}{{\text{€}}}
\providecommand{\no}{{\text{n°}}}
\providecommand{\interval}[2]{{[#1;#2]}}
\providecommand{\intervFO}[2]{{[#1;#2[}}
\providecommand{\dotfill}{\ldots}

\providecommand{\ds}{\displaystyle}
\providecommand{\textding}[1]{#1}
\providecommand{\Oij}{\(\mathrm{O};\vec{\imath},\vec{\jmath}\)}
\providecommand{\Oijk}{\(\mathrm{O};\vec{\imath},\vec{\jmath},\vec{k}\)}
\providecommand{\Ouv}{\(\mathrm{O};\vec{u},\vec{v}\)}
\providecommand{\intervFF}[2]{\ensuremath{[#1;#2]}}
\renewcommand{\euro}{\ifmmode\text{\texteuro}\else\texteuro\fi}
\DeclareUnicodeCharacter{00B0}{\ifmmode^{\circ}\else\textdegree\fi}
\DeclareUnicodeCharacter{20AC}{\ifmmode\text{\texteuro}\else\texteuro\fi}
% Énoncés : Éducation nationale. Transcription et corrigés : APMEP, Denis Vergès et contributeurs. Découpage et classement : CoopMaths, Sébastien Lozano, Éric Elter et contributeurs. Adaptation de la présentation : Maths Explorer.
% CC BY-SA : https://coopmaths.fr/www/about/
% https://coopmaths.fr/alea/static/e3c/2020/tex/e3c_2020_00_sujet49_4.tex
\begin{document}
\noindent\textit{APMEP / CoopMaths. CC BY-SA. Présentation adaptée par Maths Explorer.}\par\medskip
\medskip

On considère la fonction $f$ définie et dérivable sur $\R$ par 

\[f(x) = (ax + b)\e^{-0,1x}\]

où $a$ et $b$ sont des réels fixés.

La courbe représentative $\mathcal{C}_f$ de la fonction $f$ est donnée ci-dessous, dans un repère orthogonal.

\begin{figure}[b]
\begin{center}
\psset{xunit=0.3cm,yunit=0.3cm,labelFontSize=\scriptstyle,showorigin=false}
\begin{pspicture*}(-5,-14)(43,24)
\multido{\n=-4+2}{24}{\psline[linewidth=0.35pt,linecolor=lightgray](\n,-12)(\n,23.4)}
\multido{\n=-12+2}{18}{\psline[linewidth=0.35pt,linecolor=lightgray](-5.8,\n)(43,\n)}
\psaxes[linewidth=0.95pt,Dx=10,Dy=10]{->}(0,0)(-5,-12.5)(43.2,24)
\def\Func{x 4 mul 5 add 2.71828 x 0.1 neg mul  exp mul }
\psplot[plotpoints=2000,linewidth=1.25pt,linecolor=red]{-3.5}{43}{\Func}
\psdots[dotstyle=Mul,dotscale=1.9,linecolor=blue](0,5)(4,19)
\uput[ul](0,5){A}\uput[dl](0,0){O}\uput[ur](4,19){B}
\psplot[plotpoints=2000,linewidth=1.25pt,linecolor=blue]{-4}{5.8}{x 3.5 mul 5 add}
\rput(6,22.5){$\mathcal{T}$}
\uput[u](25,8.3){\red $\mathcal{C}_f$}
\end{pspicture*}
\end{center}
\end{figure}

On a également représenté la tangente $\mathcal{T}$ à $\mathcal{C}_f$ au point A(0~;~5).

On admet que cette tangente $\mathcal{T}$ passe par le point B(4~;~19).
\begin{enumerate}
\item  En exprimant $f(0)$, déterminer la valeur de $b$.
\item 
		\begin{enumerate}
		\item  À l'aide des coordonnées des points A et B, déterminer une équation de la droite $\mathcal{T}$.
		\item Exprimer, pour tout réel $x$, $f'(x)$ en fonction de $x$ et de $a$ et en déduire que pour tout réel $x$, $f(x)=(4x+5)\e^{-0,1x}$.
		\end{enumerate}
\item On souhaite déterminer le maximum de la fonction $f$ sur $\R$.
	\begin{enumerate}
		\item  Montrer que pour tout $x\in \R,\  f'(x)=(- 0,4x + 3,5)\e^{-0,1x}$.
		\item Déterminer les variations de $f$ sur $\R$ et en déduire le maximum de $f$ sur $\R$. 
	\end{enumerate}
\end{enumerate}

\end{document}
